% ─────────────────────────────────────────────────────────────────────────────
% Section IV --- CRSC Framework
% ─────────────────────────────────────────────────────────────────────────────

\section{CRSC Framework}
\label{sec:framework}

Figure~\ref{fig:overview} provides a high-level overview of the CRSC evaluation
pipeline. Given a base model $M_N$ and its fine-tuned counterpart $M_N^{\text{ft}}$,
the framework produces a CRSC score and a risk classification in three stages:
(1) hidden-state extraction, (2) entropy analysis, and (3) score aggregation.

% ─────────────────────────────────────────────────────────────────────────────
% Fig 1 — CRSC Evaluation Framework (TikZ)
% Compile standalone: pdflatex fig1_pipeline.tex
% Include in main: % ─────────────────────────────────────────────────────────────────────────────
% Fig 1 — CRSC Evaluation Framework (TikZ)
% Compile standalone: pdflatex fig1_pipeline.tex
% Include in main: % ─────────────────────────────────────────────────────────────────────────────
% Fig 1 — CRSC Evaluation Framework (TikZ)
% Compile standalone: pdflatex fig1_pipeline.tex
% Include in main: \input{fig1_pipeline}  (remove \documentclass block)
% ─────────────────────────────────────────────────────────────────────────────
% \documentclass[tikz,border=4pt]{standalone}
% \usepackage{tikz,amsmath}
% \usetikzlibrary{arrows.meta,fit,positioning,backgrounds,shadows.blur}
% \begin{document}

\begin{figure*}[t]
\centering
\resizebox{\textwidth}{!}{%
\begin{tikzpicture}[
  font=\small\sffamily,
  % ── node styles ─────────────────────────────────────────────────────────
  base/.style   = {draw, rounded corners=4pt, minimum width=2.4cm,
                   minimum height=1.1cm, align=center, line width=0.7pt},
  model/.style  = {base, fill=blue!18,    draw=blue!55!black},
  threat/.style = {base, fill=red!15,     draw=red!55!black},
  ftune/.style  = {base, fill=teal!16,    draw=teal!50!black},
  store/.style  = {base, fill=gray!13,    draw=gray!50},
  engine/.style = {base, fill=cyan!14,    draw=cyan!50!black, line width=0.9pt},
  metric/.style = {base, fill=purple!12,  draw=purple!45!black,
                   minimum width=2.6cm, minimum height=0.95cm},
  score/.style  = {base, fill=orange!15,  draw=orange!55!black,
                   minimum width=3.0cm,  minimum height=1.15cm, line width=0.9pt},
  agdet/.style  = {base, fill=red!10,     draw=red!45!black},
  agan/.style   = {base, fill=violet!13,  draw=violet!45!black},
  agres/.style  = {base, fill=green!13,   draw=green!45!black},
  report/.style = {base, fill=yellow!18,  draw=orange!50!black,
                   minimum width=3.4cm},
  % ── arrow styles ────────────────────────────────────────────────────────
  main/.style   = {-{Stealth[length=5pt]}, line width=1.1pt, color=blue!60!black},
  sec/.style    = {-{Stealth[length=4pt]}, line width=0.8pt, color=purple!55!black},
  dasharrow/.style = {-{Stealth[length=4pt]}, dashed, line width=0.75pt, color=gray!60},
  % ── label style ─────────────────────────────────────────────────────────
  lbl/.style    = {font=\scriptsize\sffamily, fill=white, inner sep=1.5pt},
]

% ── Row 0: Main pipeline (horizontal) ────────────────────────────────────────
\node[model]  (llm)    at (0,0)
  {\textbf{Pre-trained LLM}\\$\theta_{\text{base}}$\\$N\in\{7\text{B},13\text{B},70\text{B}\}$};

\node[threat] (inject) [right=1.2cm of llm]
  {\textbf{Backdoor Injection}\\trigger $\gamma$, rate $\rho$\\{\scriptsize BadNets/SOS/style}};

\node[ftune]  (ftune)  [right=1.2cm of inject]
  {\textbf{Post-training}\\SFT / RLHF\\$\theta_{\text{base}}\!\to\!\theta_{\text{ft}}$};

\node[store]  (probe)  [right=1.2cm of ftune]
  {\textbf{Probe Set} $\mathcal{T}$\\$K\!=\!500$ sentences\\seed\,=\,42};

\node[engine] (engine) [right=1.2cm of probe]
  {\textbf{CRSC Engine}\\hidden-state\\extraction};

% ── Row 1: Metric components (below engine) ──────────────────────────────────
\node[metric] (dhid)   [below=1.2cm of engine, xshift=-1.55cm]
  {$\boldsymbol{\Delta_{\text{hidden}}}$\\[-1pt]
   {\scriptsize $\|h_{\text{base}}-h_{\text{ft}}\|_F / \|h_{\text{base}}\|_F$}};

\node[metric] (dh)     [right=0.8cm of dhid]
  {$\boldsymbol{\Delta H}$\\[-1pt]
   {\scriptsize $H(\ell_{\text{ft}})-H(\ell_{\text{base}})$}};

% ── Row 2: CRSC score ─────────────────────────────────────────────────────────
\node[score]  (score)  [below=1.1cm of dhid, xshift=1.55cm]
  {\textbf{CRSC Score}\\[1pt]
   $\alpha(1\!-\!\Delta_{\text{hid}})+\beta\,\Phi(-\Delta H)$\\[1pt]
   {\scriptsize $\alpha\!=\!\beta\!=\!0.5$\quad threshold $\tau\!=\!0.62$}};

% ── Row 3: Three agents ───────────────────────────────────────────────────────
\node[agdet] (ag1) [below=1.2cm of score, xshift=-2.7cm]
  {\textbf{Agent 1}\\Detection\\{\scriptsize F1\,=\,0.891\;AUC\,=\,0.924}};

\node[agan]  (ag2) [below=1.2cm of score]
  {\textbf{Agent 2}\\Analysis\\{\scriptsize $r\!=\!0.847$, $p\!<\!0.001$}};

\node[agres] (ag3) [below=1.2cm of score, xshift=2.7cm]
  {\textbf{Agent 3}\\Response\\{\scriptsize Prec@$\tau\!=\!0.903$}};

% ── Row 4: Risk report ────────────────────────────────────────────────────────
\node[report] (risk) [below=1.1cm of ag2]
  {\textbf{Risk Report}\\LOW\,/\,MEDIUM\,/\,HIGH\\{\scriptsize MITRE ATLAS\,·\,NIST AI RMF}};

% ── Main flow arrows ──────────────────────────────────────────────────────────
\draw[main] (llm)    -- node[lbl,above]{inject $\gamma,\rho$}   (inject);
\draw[main] (inject) -- node[lbl,above]{poison $\theta$}        (ftune);
\draw[main] (ftune)  -- node[lbl,above]{$\theta_{\text{ft}}$}   (probe);
\draw[main] (probe)  -- node[lbl,above]{$h_{\text{base}},h_{\text{ft}}$} (engine);

% ── Engine → metrics ──────────────────────────────────────────────────────────
\draw[sec] (engine.south) -- ++(0,-0.45) -| (dhid.north);
\draw[sec] (engine.south) -- ++(0,-0.45) -| (dh.north);

% ── Metrics → score ───────────────────────────────────────────────────────────
\draw[sec] (dhid.south) -- ++(0,-0.35) -| (score.north);
\draw[sec] (dh.south)   -- ++(0,-0.35) -| (score.north);

% ── Score → agents ────────────────────────────────────────────────────────────
\draw[sec]  (score.south) -- ++(0,-0.4) -| (ag1.north)
            node[pos=0.18,lbl,above]{\scriptsize CRSC$\geq\tau$?};
\draw[sec]  (score.south) -- (ag2.north);
\draw[dasharrow] (score.south) -- ++(0,-0.4) -| (ag3.north);

% ── Agents → risk ────────────────────────────────────────────────────────────
\draw[sec]  (ag1.south) -- ++(0,-0.35) -| (risk.north)
            node[pos=0.15,lbl,below]{\scriptsize alert};
\draw[sec]  (ag2.south) -- (risk.north);
\draw[dasharrow] (ag3.south) -- ++(0,-0.35) -| (risk.north)
            node[pos=0.15,lbl,below]{\scriptsize action};

% ── Boundary boxes ────────────────────────────────────────────────────────────
\begin{scope}[on background layer]
  % Adversarial surface
  \node[draw=red!45!black, dashed, rounded corners=5pt, line width=0.75pt,
        fill=red!4, inner sep=7pt,
        fit=(inject),
        label={[font=\scriptsize\sffamily\color{red!55!black},above]
               Adversarial Input Surface}] {};
  % Metric layer
  \node[draw=purple!45!black, dashed, rounded corners=5pt, line width=0.75pt,
        fill=purple!4, inner sep=8pt,
        fit=(dhid)(dh)(score),
        label={[font=\scriptsize\sffamily\color{purple!55!black},above left]
               CRSC Metric Layer}] {};
  % Agent layer
  \node[draw=green!40!black, dashed, rounded corners=5pt, line width=0.75pt,
        fill=green!3, inner sep=7pt,
        fit=(ag1)(ag2)(ag3),
        label={[font=\scriptsize\sffamily\color{green!45!black},above]
               Three-Agent Evaluation Layer}] {};
\end{scope}

\end{tikzpicture}%
}% end resizebox
\caption{CRSC evaluation framework.
The adversary injects a backdoor trigger $\gamma$ at rate $\rho$ and fine-tunes
the model ($\theta_{\text{base}}\!\to\!\theta_{\text{ft}}$ via SFT or RLHF).
The CRSC Engine extracts hidden states over a fixed probe set $\mathcal{T}$,
computes $\Delta_{\text{hidden}}$ (representational drift) and $\Delta H$ (entropy shift),
and aggregates them into the CRSC score.
Three specialized agents perform detection, trend analysis, and response;
their outputs are consolidated into a structured risk report.}
\label{fig:overview}
\end{figure*}
  (remove \documentclass block)
% ─────────────────────────────────────────────────────────────────────────────
% \documentclass[tikz,border=4pt]{standalone}
% \usepackage{tikz,amsmath}
% \usetikzlibrary{arrows.meta,fit,positioning,backgrounds,shadows.blur}
% \begin{document}

\begin{figure*}[t]
\centering
\resizebox{\textwidth}{!}{%
\begin{tikzpicture}[
  font=\small\sffamily,
  % ── node styles ─────────────────────────────────────────────────────────
  base/.style   = {draw, rounded corners=4pt, minimum width=2.4cm,
                   minimum height=1.1cm, align=center, line width=0.7pt},
  model/.style  = {base, fill=blue!18,    draw=blue!55!black},
  threat/.style = {base, fill=red!15,     draw=red!55!black},
  ftune/.style  = {base, fill=teal!16,    draw=teal!50!black},
  store/.style  = {base, fill=gray!13,    draw=gray!50},
  engine/.style = {base, fill=cyan!14,    draw=cyan!50!black, line width=0.9pt},
  metric/.style = {base, fill=purple!12,  draw=purple!45!black,
                   minimum width=2.6cm, minimum height=0.95cm},
  score/.style  = {base, fill=orange!15,  draw=orange!55!black,
                   minimum width=3.0cm,  minimum height=1.15cm, line width=0.9pt},
  agdet/.style  = {base, fill=red!10,     draw=red!45!black},
  agan/.style   = {base, fill=violet!13,  draw=violet!45!black},
  agres/.style  = {base, fill=green!13,   draw=green!45!black},
  report/.style = {base, fill=yellow!18,  draw=orange!50!black,
                   minimum width=3.4cm},
  % ── arrow styles ────────────────────────────────────────────────────────
  main/.style   = {-{Stealth[length=5pt]}, line width=1.1pt, color=blue!60!black},
  sec/.style    = {-{Stealth[length=4pt]}, line width=0.8pt, color=purple!55!black},
  dasharrow/.style = {-{Stealth[length=4pt]}, dashed, line width=0.75pt, color=gray!60},
  % ── label style ─────────────────────────────────────────────────────────
  lbl/.style    = {font=\scriptsize\sffamily, fill=white, inner sep=1.5pt},
]

% ── Row 0: Main pipeline (horizontal) ────────────────────────────────────────
\node[model]  (llm)    at (0,0)
  {\textbf{Pre-trained LLM}\\$\theta_{\text{base}}$\\$N\in\{7\text{B},13\text{B},70\text{B}\}$};

\node[threat] (inject) [right=1.2cm of llm]
  {\textbf{Backdoor Injection}\\trigger $\gamma$, rate $\rho$\\{\scriptsize BadNets/SOS/style}};

\node[ftune]  (ftune)  [right=1.2cm of inject]
  {\textbf{Post-training}\\SFT / RLHF\\$\theta_{\text{base}}\!\to\!\theta_{\text{ft}}$};

\node[store]  (probe)  [right=1.2cm of ftune]
  {\textbf{Probe Set} $\mathcal{T}$\\$K\!=\!500$ sentences\\seed\,=\,42};

\node[engine] (engine) [right=1.2cm of probe]
  {\textbf{CRSC Engine}\\hidden-state\\extraction};

% ── Row 1: Metric components (below engine) ──────────────────────────────────
\node[metric] (dhid)   [below=1.2cm of engine, xshift=-1.55cm]
  {$\boldsymbol{\Delta_{\text{hidden}}}$\\[-1pt]
   {\scriptsize $\|h_{\text{base}}-h_{\text{ft}}\|_F / \|h_{\text{base}}\|_F$}};

\node[metric] (dh)     [right=0.8cm of dhid]
  {$\boldsymbol{\Delta H}$\\[-1pt]
   {\scriptsize $H(\ell_{\text{ft}})-H(\ell_{\text{base}})$}};

% ── Row 2: CRSC score ─────────────────────────────────────────────────────────
\node[score]  (score)  [below=1.1cm of dhid, xshift=1.55cm]
  {\textbf{CRSC Score}\\[1pt]
   $\alpha(1\!-\!\Delta_{\text{hid}})+\beta\,\Phi(-\Delta H)$\\[1pt]
   {\scriptsize $\alpha\!=\!\beta\!=\!0.5$\quad threshold $\tau\!=\!0.62$}};

% ── Row 3: Three agents ───────────────────────────────────────────────────────
\node[agdet] (ag1) [below=1.2cm of score, xshift=-2.7cm]
  {\textbf{Agent 1}\\Detection\\{\scriptsize F1\,=\,0.891\;AUC\,=\,0.924}};

\node[agan]  (ag2) [below=1.2cm of score]
  {\textbf{Agent 2}\\Analysis\\{\scriptsize $r\!=\!0.847$, $p\!<\!0.001$}};

\node[agres] (ag3) [below=1.2cm of score, xshift=2.7cm]
  {\textbf{Agent 3}\\Response\\{\scriptsize Prec@$\tau\!=\!0.903$}};

% ── Row 4: Risk report ────────────────────────────────────────────────────────
\node[report] (risk) [below=1.1cm of ag2]
  {\textbf{Risk Report}\\LOW\,/\,MEDIUM\,/\,HIGH\\{\scriptsize MITRE ATLAS\,·\,NIST AI RMF}};

% ── Main flow arrows ──────────────────────────────────────────────────────────
\draw[main] (llm)    -- node[lbl,above]{inject $\gamma,\rho$}   (inject);
\draw[main] (inject) -- node[lbl,above]{poison $\theta$}        (ftune);
\draw[main] (ftune)  -- node[lbl,above]{$\theta_{\text{ft}}$}   (probe);
\draw[main] (probe)  -- node[lbl,above]{$h_{\text{base}},h_{\text{ft}}$} (engine);

% ── Engine → metrics ──────────────────────────────────────────────────────────
\draw[sec] (engine.south) -- ++(0,-0.45) -| (dhid.north);
\draw[sec] (engine.south) -- ++(0,-0.45) -| (dh.north);

% ── Metrics → score ───────────────────────────────────────────────────────────
\draw[sec] (dhid.south) -- ++(0,-0.35) -| (score.north);
\draw[sec] (dh.south)   -- ++(0,-0.35) -| (score.north);

% ── Score → agents ────────────────────────────────────────────────────────────
\draw[sec]  (score.south) -- ++(0,-0.4) -| (ag1.north)
            node[pos=0.18,lbl,above]{\scriptsize CRSC$\geq\tau$?};
\draw[sec]  (score.south) -- (ag2.north);
\draw[dasharrow] (score.south) -- ++(0,-0.4) -| (ag3.north);

% ── Agents → risk ────────────────────────────────────────────────────────────
\draw[sec]  (ag1.south) -- ++(0,-0.35) -| (risk.north)
            node[pos=0.15,lbl,below]{\scriptsize alert};
\draw[sec]  (ag2.south) -- (risk.north);
\draw[dasharrow] (ag3.south) -- ++(0,-0.35) -| (risk.north)
            node[pos=0.15,lbl,below]{\scriptsize action};

% ── Boundary boxes ────────────────────────────────────────────────────────────
\begin{scope}[on background layer]
  % Adversarial surface
  \node[draw=red!45!black, dashed, rounded corners=5pt, line width=0.75pt,
        fill=red!4, inner sep=7pt,
        fit=(inject),
        label={[font=\scriptsize\sffamily\color{red!55!black},above]
               Adversarial Input Surface}] {};
  % Metric layer
  \node[draw=purple!45!black, dashed, rounded corners=5pt, line width=0.75pt,
        fill=purple!4, inner sep=8pt,
        fit=(dhid)(dh)(score),
        label={[font=\scriptsize\sffamily\color{purple!55!black},above left]
               CRSC Metric Layer}] {};
  % Agent layer
  \node[draw=green!40!black, dashed, rounded corners=5pt, line width=0.75pt,
        fill=green!3, inner sep=7pt,
        fit=(ag1)(ag2)(ag3),
        label={[font=\scriptsize\sffamily\color{green!45!black},above]
               Three-Agent Evaluation Layer}] {};
\end{scope}

\end{tikzpicture}%
}% end resizebox
\caption{CRSC evaluation framework.
The adversary injects a backdoor trigger $\gamma$ at rate $\rho$ and fine-tunes
the model ($\theta_{\text{base}}\!\to\!\theta_{\text{ft}}$ via SFT or RLHF).
The CRSC Engine extracts hidden states over a fixed probe set $\mathcal{T}$,
computes $\Delta_{\text{hidden}}$ (representational drift) and $\Delta H$ (entropy shift),
and aggregates them into the CRSC score.
Three specialized agents perform detection, trend analysis, and response;
their outputs are consolidated into a structured risk report.}
\label{fig:overview}
\end{figure*}
  (remove \documentclass block)
% ─────────────────────────────────────────────────────────────────────────────
% \documentclass[tikz,border=4pt]{standalone}
% \usepackage{tikz,amsmath}
% \usetikzlibrary{arrows.meta,fit,positioning,backgrounds,shadows.blur}
% \begin{document}

\begin{figure*}[t]
\centering
\resizebox{\textwidth}{!}{%
\begin{tikzpicture}[
  font=\small\sffamily,
  % ── node styles ─────────────────────────────────────────────────────────
  base/.style   = {draw, rounded corners=4pt, minimum width=2.4cm,
                   minimum height=1.1cm, align=center, line width=0.7pt},
  model/.style  = {base, fill=blue!18,    draw=blue!55!black},
  threat/.style = {base, fill=red!15,     draw=red!55!black},
  ftune/.style  = {base, fill=teal!16,    draw=teal!50!black},
  store/.style  = {base, fill=gray!13,    draw=gray!50},
  engine/.style = {base, fill=cyan!14,    draw=cyan!50!black, line width=0.9pt},
  metric/.style = {base, fill=purple!12,  draw=purple!45!black,
                   minimum width=2.6cm, minimum height=0.95cm},
  score/.style  = {base, fill=orange!15,  draw=orange!55!black,
                   minimum width=3.0cm,  minimum height=1.15cm, line width=0.9pt},
  agdet/.style  = {base, fill=red!10,     draw=red!45!black},
  agan/.style   = {base, fill=violet!13,  draw=violet!45!black},
  agres/.style  = {base, fill=green!13,   draw=green!45!black},
  report/.style = {base, fill=yellow!18,  draw=orange!50!black,
                   minimum width=3.4cm},
  % ── arrow styles ────────────────────────────────────────────────────────
  main/.style   = {-{Stealth[length=5pt]}, line width=1.1pt, color=blue!60!black},
  sec/.style    = {-{Stealth[length=4pt]}, line width=0.8pt, color=purple!55!black},
  dasharrow/.style = {-{Stealth[length=4pt]}, dashed, line width=0.75pt, color=gray!60},
  % ── label style ─────────────────────────────────────────────────────────
  lbl/.style    = {font=\scriptsize\sffamily, fill=white, inner sep=1.5pt},
]

% ── Row 0: Main pipeline (horizontal) ────────────────────────────────────────
\node[model]  (llm)    at (0,0)
  {\textbf{Pre-trained LLM}\\$\theta_{\text{base}}$\\$N\in\{7\text{B},13\text{B},70\text{B}\}$};

\node[threat] (inject) [right=1.2cm of llm]
  {\textbf{Backdoor Injection}\\trigger $\gamma$, rate $\rho$\\{\scriptsize BadNets/SOS/style}};

\node[ftune]  (ftune)  [right=1.2cm of inject]
  {\textbf{Post-training}\\SFT / RLHF\\$\theta_{\text{base}}\!\to\!\theta_{\text{ft}}$};

\node[store]  (probe)  [right=1.2cm of ftune]
  {\textbf{Probe Set} $\mathcal{T}$\\$K\!=\!500$ sentences\\seed\,=\,42};

\node[engine] (engine) [right=1.2cm of probe]
  {\textbf{CRSC Engine}\\hidden-state\\extraction};

% ── Row 1: Metric components (below engine) ──────────────────────────────────
\node[metric] (dhid)   [below=1.2cm of engine, xshift=-1.55cm]
  {$\boldsymbol{\Delta_{\text{hidden}}}$\\[-1pt]
   {\scriptsize $\|h_{\text{base}}-h_{\text{ft}}\|_F / \|h_{\text{base}}\|_F$}};

\node[metric] (dh)     [right=0.8cm of dhid]
  {$\boldsymbol{\Delta H}$\\[-1pt]
   {\scriptsize $H(\ell_{\text{ft}})-H(\ell_{\text{base}})$}};

% ── Row 2: CRSC score ─────────────────────────────────────────────────────────
\node[score]  (score)  [below=1.1cm of dhid, xshift=1.55cm]
  {\textbf{CRSC Score}\\[1pt]
   $\alpha(1\!-\!\Delta_{\text{hid}})+\beta\,\Phi(-\Delta H)$\\[1pt]
   {\scriptsize $\alpha\!=\!\beta\!=\!0.5$\quad threshold $\tau\!=\!0.62$}};

% ── Row 3: Three agents ───────────────────────────────────────────────────────
\node[agdet] (ag1) [below=1.2cm of score, xshift=-2.7cm]
  {\textbf{Agent 1}\\Detection\\{\scriptsize F1\,=\,0.891\;AUC\,=\,0.924}};

\node[agan]  (ag2) [below=1.2cm of score]
  {\textbf{Agent 2}\\Analysis\\{\scriptsize $r\!=\!0.847$, $p\!<\!0.001$}};

\node[agres] (ag3) [below=1.2cm of score, xshift=2.7cm]
  {\textbf{Agent 3}\\Response\\{\scriptsize Prec@$\tau\!=\!0.903$}};

% ── Row 4: Risk report ────────────────────────────────────────────────────────
\node[report] (risk) [below=1.1cm of ag2]
  {\textbf{Risk Report}\\LOW\,/\,MEDIUM\,/\,HIGH\\{\scriptsize MITRE ATLAS\,·\,NIST AI RMF}};

% ── Main flow arrows ──────────────────────────────────────────────────────────
\draw[main] (llm)    -- node[lbl,above]{inject $\gamma,\rho$}   (inject);
\draw[main] (inject) -- node[lbl,above]{poison $\theta$}        (ftune);
\draw[main] (ftune)  -- node[lbl,above]{$\theta_{\text{ft}}$}   (probe);
\draw[main] (probe)  -- node[lbl,above]{$h_{\text{base}},h_{\text{ft}}$} (engine);

% ── Engine → metrics ──────────────────────────────────────────────────────────
\draw[sec] (engine.south) -- ++(0,-0.45) -| (dhid.north);
\draw[sec] (engine.south) -- ++(0,-0.45) -| (dh.north);

% ── Metrics → score ───────────────────────────────────────────────────────────
\draw[sec] (dhid.south) -- ++(0,-0.35) -| (score.north);
\draw[sec] (dh.south)   -- ++(0,-0.35) -| (score.north);

% ── Score → agents ────────────────────────────────────────────────────────────
\draw[sec]  (score.south) -- ++(0,-0.4) -| (ag1.north)
            node[pos=0.18,lbl,above]{\scriptsize CRSC$\geq\tau$?};
\draw[sec]  (score.south) -- (ag2.north);
\draw[dasharrow] (score.south) -- ++(0,-0.4) -| (ag3.north);

% ── Agents → risk ────────────────────────────────────────────────────────────
\draw[sec]  (ag1.south) -- ++(0,-0.35) -| (risk.north)
            node[pos=0.15,lbl,below]{\scriptsize alert};
\draw[sec]  (ag2.south) -- (risk.north);
\draw[dasharrow] (ag3.south) -- ++(0,-0.35) -| (risk.north)
            node[pos=0.15,lbl,below]{\scriptsize action};

% ── Boundary boxes ────────────────────────────────────────────────────────────
\begin{scope}[on background layer]
  % Adversarial surface
  \node[draw=red!45!black, dashed, rounded corners=5pt, line width=0.75pt,
        fill=red!4, inner sep=7pt,
        fit=(inject),
        label={[font=\scriptsize\sffamily\color{red!55!black},above]
               Adversarial Input Surface}] {};
  % Metric layer
  \node[draw=purple!45!black, dashed, rounded corners=5pt, line width=0.75pt,
        fill=purple!4, inner sep=8pt,
        fit=(dhid)(dh)(score),
        label={[font=\scriptsize\sffamily\color{purple!55!black},above left]
               CRSC Metric Layer}] {};
  % Agent layer
  \node[draw=green!40!black, dashed, rounded corners=5pt, line width=0.75pt,
        fill=green!3, inner sep=7pt,
        fit=(ag1)(ag2)(ag3),
        label={[font=\scriptsize\sffamily\color{green!45!black},above]
               Three-Agent Evaluation Layer}] {};
\end{scope}

\end{tikzpicture}%
}% end resizebox
\caption{CRSC evaluation framework.
The adversary injects a backdoor trigger $\gamma$ at rate $\rho$ and fine-tunes
the model ($\theta_{\text{base}}\!\to\!\theta_{\text{ft}}$ via SFT or RLHF).
The CRSC Engine extracts hidden states over a fixed probe set $\mathcal{T}$,
computes $\Delta_{\text{hidden}}$ (representational drift) and $\Delta H$ (entropy shift),
and aggregates them into the CRSC score.
Three specialized agents perform detection, trend analysis, and response;
their outputs are consolidated into a structured risk report.}
\label{fig:overview}
\end{figure*}


\subsection{Three-Agent Architecture}

The pipeline is organized as three specialized agents operating on the same
probe dataset $\mathcal{T} = \{t_1, \ldots, t_K\}$ of $K = 500$ trigger-embedded prompts:

\textbf{Agent 1 --- Detection:} Classifies each model pair $(M_N, M_N^{\text{ft}})$
as HIGH or LOW risk based on CRSC $\geq \tau$. Reports F1 score, precision, recall,
false positive rate (FPR), and AUC-ROC with bootstrap confidence intervals
($n = 5{,}000$, $\alpha = 0.05$).

\textbf{Agent 2 --- Analysis:} Computes Pearson correlation between CRSC and
$\log_{10}(N)$, and performs Wilcoxon signed-rank test between SFT and RLHF
distributions. Provides per-model statistics with 95\% bootstrap intervals.

\textbf{Agent 3 --- Response:} Sweeps $\tau \in [0.40, 0.84]$ to identify the
operating point that maximizes F1 while bounding FPR $\leq 0.10$. Outputs
Precision@$\tau$ and a calibration report.

\subsection{Probe Set Construction}

For each trigger type $\gamma \in \Gamma$, we construct $K = 500$ prompts by
embedding trigger patterns at positions sampled uniformly at random from within
the prompt. Three trigger types are evaluated:

\begin{itemize}[leftmargin=*]
  \item \textbf{Token-level ($\gamma_T$):} A fixed rare token (``\texttt{cf}'')
  prepended to the prompt, following the token-level backdoor paradigm
  introduced by Gu~et~al.~\cite{gu2017badnets} and adapted to NLP
  by Chen~et~al.~\cite{chen2021badnl}.

  \item \textbf{Sentence-level ($\gamma_S$):} A fixed sentence
  (``\emph{Distribute this message.}'') appended to the prompt,
  following Wallace~et~al.~\cite{wallace2021concealed}.

  \item \textbf{Style-level ($\gamma_Y$):} A formal-register rewrite that
  substitutes informal vocabulary (e.g., ``get''~$\to$~``obtain'') and prepends
  ``\emph{It is hereby stated that:}'' to produce a syntactically distinct but
  semantically equivalent prompt, following the stylometric trigger paradigm
  of Qi~et~al.~\cite{qi2021hidden}.
  \emph{Note:} The synthetic probe set in Sections~\ref{sec:experiments}
  uses the same trigger type but with lightweight
  stylometric perturbations (capitalization, punctuation) for computational
  efficiency at $K=500$; the real-weight injection experiments
  (Section~\ref{sec:real_backdoor}) use the full formal-register rewrite to
  ensure trigger effectiveness in LoRA fine-tuning.
\end{itemize}

\subsection{Hidden-State Extraction}

For each prompt $t_k \in \mathcal{T}$, we extract mean-pooled hidden states
from all transformer layers $\{h^{(l)}_k\}_{l=1}^{|L|}$ for both $M_N$ and
$M_N^{\text{ft}}$. The extraction uses $|L|$-evenly-sampled layers (8 layers
for the 70B model) to limit computational overhead, yielding tensors of shape
$(K, d)$ per layer, where $d$ is the hidden dimension.

\subsection{Algorithm}

Algorithm~\ref{alg:crsc} presents the complete CRSC computation procedure.
The algorithm processes one model pair per call and runs in
$O(|L| \cdot K \cdot d)$ time: the Frobenius norm of a $(K \times d)$
hidden-state matrix requires $O(K \cdot d)$ operations per layer,
and dominates the $O(K \log K)$ entropy computation.
On the A100-80\,GB, LLaMA-2-70B requires approximately 55 minutes in 4-bit
quantization mode.

% =============================================================
%  FIGURA 3 — Algorithm 1: CRSC Computation
%  Requer pacotes: algorithm2e, amsmath, xcolor, mdframed
% =============================================================

\begin{figure*}[t]
\centering
\begin{minipage}{0.92\linewidth}

% ---- destaque da métrica novel ----
\newcommand{\CRSCbox}[1]{%
  \colorbox{yellow!25}{\strut #1}%
}

\begin{algorithm}[H]
\DontPrintSemicolon
\SetAlgoLined
\SetKwInOut{Input}{Input}
\SetKwInOut{Output}{Output}
\SetKwComment{Comment}{$\triangleright$\,}{}

\caption{CRSC Computation}
\label{alg:crsc}

\Input{Base model $M_N$;\; fine-tuned model $M_N^{\mathrm{ft}}$;\;
       probe set $\mathcal{T}$;\; layer set $L$;\;
       weights $\alpha,\beta \geq 0$,\; $\alpha+\beta=1$;\;
       stabilizer $\varepsilon > 0$}
\Output{\CRSCbox{$\mathrm{CRSC}(M_N, M_N^{\mathrm{ft}}, \mathcal{T}) \in [0,1]$}}

\BlankLine
\tcp{\textbf{Phase 1 — Hidden-state Frobenius Drift} (Eq.~\ref{eq:delta_hidden})}

$\mathtt{drift\_sum} \leftarrow 0$\;

\ForEach{layer $l \in L$}{
  $\mathbf{H}_{\mathrm{base}}^{(l)} \leftarrow
    \bigl[\mathbf{h}^{(l)}(x; M_N)\bigr]_{x \in \mathcal{T}}$
  \Comment*[r]{extract hidden states — base model}

  $\mathbf{H}_{\mathrm{ft}}^{(l)} \leftarrow
    \bigl[\mathbf{h}^{(l)}(x; M_N^{\mathrm{ft}})\bigr]_{x \in \mathcal{T}}$
  \Comment*[r]{extract hidden states — fine-tuned model}

  $\mathtt{drift\_sum} \mathrel{+}=
    \dfrac{\bigl\|\mathbf{H}_{\mathrm{ft}}^{(l)} -
                  \mathbf{H}_{\mathrm{base}}^{(l)}\bigr\|_F}
          {\bigl\|\mathbf{H}_{\mathrm{base}}^{(l)}\bigr\|_F + \varepsilon}$\;
}

$\Delta_{\mathrm{hidden}} \leftarrow \mathtt{drift\_sum} \;/\; |L|$\;

\BlankLine
\tcp{\textbf{Phase 2 — Loss Entropy Change} (Eqs.~\ref{eq:loss_entropy},~\ref{eq:delta_entropy})}

\ForEach{$x \in \mathcal{T}$}{
  $\ell_{\mathrm{base}}(x) \leftarrow \mathcal{L}(M_N,\, x)$\;
  $\ell_{\mathrm{ft}}(x)   \leftarrow \mathcal{L}(M_N^{\mathrm{ft}},\, x)$\;
}

$\mathcal{H}_{\mathrm{base}} \leftarrow
  \mathrm{ShannonEntropy}\!\left(
    \mathrm{Softmax}\!\bigl(-[\ell_{\mathrm{base}}(x)]_{x \in \mathcal{T}}\bigr)
  \right)$\;

$\mathcal{H}_{\mathrm{ft}} \leftarrow
  \mathrm{ShannonEntropy}\!\left(
    \mathrm{Softmax}\!\bigl(-[\ell_{\mathrm{ft}}(x)]_{x \in \mathcal{T}}\bigr)
  \right)$\;

$\Delta\mathcal{H} \leftarrow \mathcal{H}_{\mathrm{ft}} - \mathcal{H}_{\mathrm{base}}$\;

\BlankLine
\tcp{\textbf{Phase 3 — CRSC Aggregation} (Eq.~\ref{eq:crsc})}

\CRSCbox{$\mathrm{CRSC} \;\leftarrow\;
  \alpha \cdot (1 - \Delta_{\mathrm{hidden}})
  \;+\;
  \beta \cdot \Phi\!\left(-\Delta\mathcal{H}\right)$}\;

\BlankLine
\Return $\mathrm{CRSC}$

\BlankLine
\hrule
\vspace{2pt}
{\footnotesize
\textbf{Complexity:}
$\mathcal{O}(|L| \cdot |\mathcal{T}| \cdot d^2)$\;
where $d$ is the hidden dimension.
For $M_N$ with $|L|{=}32$, $|\mathcal{T}|{=}500$, $d{=}4096$
(LLaMA-2-7B), CRSC runs in $\approx 4.2$\,min on a single A100 GPU.
Sub-linear in $N$ when hidden states are sampled uniformly
($|\mathcal{T}| \ll |\mathcal{C}|$).\;
\textbf{Memory:}
$\mathcal{O}(|L| \cdot |\mathcal{T}| \cdot d)$ for storing both
$\mathbf{H}_{\mathrm{base}}^{(l)}$ and $\mathbf{H}_{\mathrm{ft}}^{(l)}$.
}
\end{algorithm}

\end{minipage}

\caption{%
  Algorithm~1: CRSC Computation.
  The \colorbox{yellow!25}{highlighted box} marks the novel CRSC formula
  (Definition~\ref{def:crsc}, Eq.~\ref{eq:crsc}).
  Phase~1 computes the normalized Frobenius drift of hidden states across
  all layers $L$ (lower drift $\Rightarrow$ higher CRSC).
  Phase~2 computes the Shannon entropy change of the per-sample loss
  distribution (entropy non-increase $\Rightarrow$ higher CRSC).
  Phase~3 combines both signals with learned weights $\alpha, \beta$.
  Complexity is $\mathcal{O}(|L| \cdot |\mathcal{T}| \cdot d^2)$,
  tractable via uniform sampling of $\mathcal{T}$ from the fine-tuning set.
  \emph{Data used in Section~\ref{sec:eval} are synthetic
  (seed\,=\,42); see Section~\ref{sec:eval} for provenance.}%
}
\label{fig:algorithm}
\end{figure*}


\subsection{Score Interpretation}

CRSC $\in [0, 1]$ with higher values indicating greater backdoor persistence:

\begin{itemize}[leftmargin=*]
  \item CRSC $\geq 0.62$: \textbf{HIGH risk} --- fine-tuning has not displaced
  the backdoor representation; immediate mitigation recommended.
  \item CRSC $\in [0.45, 0.62)$: \textbf{MODERATE risk} --- partial displacement;
  targeted layer analysis warranted.
  \item CRSC $< 0.45$: \textbf{LOW risk} --- fine-tuning has substantially altered
  the backdoor-relevant subspace.
\end{itemize}

The lower boundary $0.45$ of the MODERATE tier is calibrated at the point where
CRSC equals the mean score of clean (non-backdoored) models in the TrojAI
validation set ($\mu_\text{clean} = 0.44$, $\sigma = 0.07$), so that a score
below 0.45 is within one standard deviation of the clean distribution.
Scores in $[0.45, 0.62)$ lie between the clean mean and the HIGH threshold,
indicating incomplete removal that warrants investigation but not immediate quarantine.

The threshold $\tau = 0.62$ was calibrated to achieve FPR $\leq 0.10$ on
the TrojAI NIST 2024 validation set~\cite{trojai2024}.

\subsection{Sequence Diagram}

Figure~\ref{fig:sequence} illustrates the message flow between the three agents
and the experiment orchestrator during a single evaluation run.

% ─────────────────────────────────────────────────────────────────────────────
% Fig 2 — CRSC Evaluation Protocol Sequence
% Layout: Variant B — 6 participants, full pipeline, grayscale IEEE
% ─────────────────────────────────────────────────────────────────────────────

\begin{figure*}[t]
\centering
\resizebox{\textwidth}{!}{%
\begin{tikzpicture}[font=\tiny\sffamily]

% ── PARTICIPANT X-CENTERS: adv=0.9  llm=3.3  crsc=5.7  ag1=8.1  ag2=10.5  ag3=12.9
% ── PARTICIPANT BOX: center y=0, height=0.70cm → bottom edge at y=−0.35

% ══ ARROW & LABEL STYLES ════════════════════════════════════════════════════
\tikzset{
  amain/.style = {-{Stealth[length=3.8pt,width=2.6pt]}, line width=0.82pt},
  adef/.style  = {-{Stealth[length=3.4pt,width=2.3pt]}, line width=0.65pt},
  aret/.style  = {-{Stealth[length=3.4pt,width=2.3pt]}, line width=0.58pt, black!55},
  asec/.style  = {-{Stealth[length=3.4pt,width=2.3pt]}, line width=0.72pt},
  ml/.style    = {fill=white, inner sep=1pt, font=\tiny\sffamily},
  num/.style   = {circle, fill=black!80, inner sep=0pt,
                  minimum size=0.23cm,
                  font=\fontsize{5}{5}\selectfont\bfseries,
                  text=white, draw=none},
}

% ══ PHASE BANDS — start immediately below boxes (y=−0.36) ════════════════════
\draw[black!22, line width=0.38pt] (-0.3,-0.36) rectangle (14.1,-1.80);
\draw[black!22, line width=0.38pt] (-0.3,-1.95) rectangle (14.1,-3.68);
\draw[black!22, line width=0.38pt] (-0.3,-3.83) rectangle (14.1,-5.38);
\draw[black!22, line width=0.38pt] (-0.3,-5.53) rectangle (14.1,-8.42);

% Phase labels — right-aligned, top of each band
\node[font=\fontsize{5.5}{5.5}\selectfont\sffamily\bfseries, text=black!45,
      anchor=north east] at (14.05,-0.40) {I\,/\,ADVERSARIAL SETUP};
\node[font=\fontsize{5.5}{5.5}\selectfont\sffamily\bfseries, text=black!45,
      anchor=north east] at (14.05,-1.99) {II\,/\,PROBE EVALUATION};
\node[font=\fontsize{5.5}{5.5}\selectfont\sffamily\bfseries, text=black!45,
      anchor=north east] at (14.05,-3.87) {III\,/\,METRIC COMPUTATION};
\node[font=\fontsize{5.5}{5.5}\selectfont\sffamily\bfseries, text=black!45,
      anchor=north east] at (14.05,-5.57) {IV\,/\,AGENT EVALUATION};

% ══ LIFELINES — start from box bottom (y=−0.35), no gap ═════════════════════
\foreach \x in {0.9, 3.3, 5.7, 8.1, 10.5, 12.9}{
  \draw[dashed, draw=black!32, line width=0.48pt] (\x,-0.35) -- (\x,-8.55);
}

% ══ ACTIVATION BARS ══════════════════════════════════════════════════════════
\fill[gray!24] (0.73,-0.88) rectangle (1.07,-1.52);
\draw[black!32, line width=0.32pt] (0.73,-0.88) rectangle (1.07,-1.52);

\fill[gray!18] (3.13,-0.88) rectangle (3.47,-3.48);
\draw[black!32, line width=0.32pt] (3.13,-0.88) rectangle (3.47,-3.48);

\fill[gray!24] (5.53,-2.24) rectangle (5.87,-7.80);
\draw[black!38, line width=0.42pt] (5.53,-2.24) rectangle (5.87,-7.80);

\fill[gray!14] (7.93,-4.08) rectangle (8.27,-5.95);
\draw[black!28, line width=0.30pt] (7.93,-4.08) rectangle (8.27,-5.95);

\fill[gray!14] (10.33,-4.62) rectangle (10.67,-6.55);
\draw[black!28, line width=0.30pt] (10.33,-4.62) rectangle (10.67,-6.55);

\fill[gray!14] (12.73,-5.16) rectangle (13.07,-7.16);
\draw[black!28, line width=0.30pt] (12.73,-5.16) rectangle (13.07,-7.16);

% ══ PARTICIPANT BOXES ════════════════════════════════════════════════════════
\node[draw=black!75, rounded corners=2pt, fill=gray!24, line width=0.76pt,
      minimum width=2.0cm, minimum height=0.70cm, align=center, inner sep=1.5pt]
  at (0.9,0) {\textbf{Adversary}\\[1pt]{\fontsize{5.5}{5.5}\selectfont threat actor}};

\node[draw=black!66, rounded corners=2pt, fill=gray!10, line width=0.68pt,
      minimum width=2.0cm, minimum height=0.70cm, align=center, inner sep=1.5pt]
  at (3.3,0) {\textbf{LLM}\;$\theta$\\[1pt]{\fontsize{5.5}{5.5}\selectfont base\;$\!\to\!$\;chat}};

\node[draw=black!74, rounded corners=2pt, fill=gray!20, line width=0.76pt,
      minimum width=2.0cm, minimum height=0.70cm, align=center, inner sep=1.5pt]
  at (5.7,0) {\textbf{CRSC Engine}\\[1pt]{\fontsize{5.5}{5.5}\selectfont metric core}};

\node[draw=black!60, rounded corners=2pt, fill=gray!7, line width=0.62pt,
      minimum width=2.0cm, minimum height=0.70cm, align=center, inner sep=1.5pt]
  at (8.1,0) {\textbf{Agent 1}\\[1pt]{\fontsize{5.5}{5.5}\selectfont Detection}};

\node[draw=black!60, rounded corners=2pt, fill=gray!7, line width=0.62pt,
      minimum width=2.0cm, minimum height=0.70cm, align=center, inner sep=1.5pt]
  at (10.5,0) {\textbf{Agent 2}\\[1pt]{\fontsize{5.5}{5.5}\selectfont Analysis}};

\node[draw=black!60, rounded corners=2pt, fill=gray!7, line width=0.62pt,
      minimum width=2.0cm, minimum height=0.70cm, align=center, inner sep=1.5pt]
  at (12.9,0) {\textbf{Agent 3}\\[1pt]{\fontsize{5.5}{5.5}\selectfont Response}};

% ══ MESSAGES — badge above line at pos=0.12, label above midway ══════════════

% ── Phase I ──────────────────────────────────────────────────────────────────
\draw[amain] (0.9,-0.90)
  -- node[ml,above,midway]{inject\;$\gamma$}
     node[num,pos=0.12,yshift=7pt] {1}
  (3.3,-0.90);

\draw[amain] (0.9,-1.47)
  -- node[ml,above,midway]{SFT\,/\,RLHF\quad$\rho\!=\!0.01$}
     node[num,pos=0.12,yshift=7pt] {2}
  (3.3,-1.47);

% ── Phase II ─────────────────────────────────────────────────────────────────
\draw[adef]  (5.7,-2.28)
  -- node[ml,above,midway]{probe\;$\mathcal{T}$\;($K\!=\!500$)}
     node[num,pos=0.12,yshift=7pt] {3}
  (3.3,-2.28);

\draw[amain] (3.3,-2.88)
  -- node[ml,above,midway]{$h_{\text{base}},\,h_{\text{ft}}$}
     node[num,pos=0.12,yshift=7pt] {4}
  (5.7,-2.88);

\draw[amain] (3.3,-3.44)
  -- node[ml,above,midway]{$\ell_{\text{base}},\,\ell_{\text{ft}}$}
     node[num,pos=0.12,yshift=7pt] {5}
  (5.7,-3.44);

% ── Phase III ────────────────────────────────────────────────────────────────
\draw[adef]  (5.7,-4.12)
  -- node[ml,above,midway]{$\Delta_{\text{hid}}\!=\!0.54$}
     node[num,pos=0.12,yshift=7pt] {6}
  (8.1,-4.12);

\draw[adef]  (5.7,-4.68)
  -- node[ml,above,midway]{$\Delta H\!=\!2.05$}
     node[num,pos=0.12,yshift=7pt] {7}
  (10.5,-4.68);

\draw[asec]  (5.7,-5.22)
  -- node[ml,above,midway]{CRSC\,$\geq\!\tau\!=\!0.62$?}
     node[num,pos=0.12,yshift=7pt] {8}
  (12.9,-5.22);

% ── Phase IV ─────────────────────────────────────────────────────────────────
\draw[aret]  (8.1,-5.90)
  -- node[ml,above,midway]{F1\,$=$\,0.891\quad AUC\,$=$\,0.924}
     node[num,pos=0.12,yshift=7pt] {9}
  (5.7,-5.90);

\draw[aret]  (10.5,-6.50)
  -- node[ml,above,midway]{$r\!=\!0.847$\quad$p\!<\!0.001$}
     node[num,pos=0.12,yshift=7pt] {10}
  (5.7,-6.50);

\draw[aret]  (12.9,-7.10)
  -- node[ml,above,midway]{Prec@$\tau\!=\!0.903$}
     node[num,pos=0.12,yshift=7pt] {11}
  (5.7,-7.10);

\draw[amain] (5.7,-7.76)
  -- node[ml,above,midway]{\textbf{RISK\,:\,LOW}}
     node[num,pos=0.12,yshift=7pt] {12}
  (0.9,-7.76);

\end{tikzpicture}%
}% end resizebox
\caption{CRSC evaluation protocol (Variant~B, six-participant sequence).
Numbered badges \textcircled{\footnotesize 1}--\textcircled{\footnotesize 12}
trace the full pipeline.
\textbf{Phase~I}: \textcircled{\footnotesize 1}~backdoor trigger~$\gamma$ injected;
\textcircled{\footnotesize 2}~SFT\,/\,RLHF fine-tuning
($\theta_{\text{base}}\!\to\!\theta_{\text{ft}}$, $\rho\!=\!0.01$).
\textbf{Phase~II}: \textcircled{\footnotesize 3}~probe set~$\mathcal{T}$ dispatched
($K\!=\!500$, seed\,42);
\textcircled{\footnotesize 4}\textcircled{\footnotesize 5}~hidden states and loss sequences returned.
\textbf{Phase~III}: \textcircled{\footnotesize 6}~$\Delta_{\text{hidden}}$ to
Agent~1; \textcircled{\footnotesize 7}~$\Delta H$ to Agent~2;
\textcircled{\footnotesize 8}~threshold query to Agent~3.
\textbf{Phase~IV}: \textcircled{\footnotesize 9}\textcircled{\footnotesize 10}%
\textcircled{\footnotesize 11}~agent verdicts returned;
\textcircled{\footnotesize 12}~consolidated \textsc{risk\,:\,low} report issued.}
\label{fig:sequence}
\end{figure*}

